\section{Conclusi\'{o}n}
\label{sec:conclusion}

Esta versi\'{o}n del art\'{i}culo presenta resultados preliminares para un solo pa\'{i}s. Se reportan las estimaciones causales para Per\'{u} (Lima Metropolitana, EPEN trimestral 2021T1--2025T4), construidas sobre el panel armonizado a CIUO-08 a cuatro d\'{i}gitos descrito en la \cref{sec:data}. La aplicaci\'{o}n del mismo dise\~{n}o a Costa Rica, Colombia, Ecuador, M\'{e}xico y Uruguay queda pendiente, y con ello tambi\'{e}n queda pendiente la comparaci\'{o}n entre pa\'{i}ses que constituye la contribuci\'{o}n emp\'{i}rica central del art\'{i}culo. Las afirmaciones a continuaci\'{o}n deben leerse como caracterizaci\'{o}n de Lima Metropolitana, no como evidencia regional para Am\'{e}rica Latina.

Tres conclusiones emergen de la evidencia peruana. Primero, la lectura agregada del shock de los LLMs en Lima Metropolitana es enga\~{n}osamente nula: la especificaci\'{o}n continua TWFE arroja coeficientes peque\~{n}os y estad\'{i}sticamente indistinguibles de cero al 95 por ciento sobre los tres resultados primarios, lo que en aislamiento sugerir\'{i}a una replicaci\'{o}n del nulo de \citet{HumlumVestergaard2025_llm}. La especificaci\'{o}n binaria de \citet{CallawayS2021_did} sobre la submuestra balanceada matiza esa lectura entregando un efecto positivo sobre el empleo ($+14.4$ por ciento, $p < 0.01$) y un efecto negativo sobre las horas semanales ($-0.80$, $p < 0.05$). La heterogeneidad por estimador es informativa por s\'{i} misma sobre las limitaciones de TWFE en presencia de respuestas heterog\'{e}neas a la dosis.

Segundo, la heterogeneidad por formalidad de l\'{i}nea base es la fuente principal del cuadro emp\'{i}rico. Los canales formal e informal se mueven en direcciones opuestas y estad\'{i}sticamente significativas en empleo ($\hat\tau_F = +0.567$ vs.\ $\hat\tau_I = -0.222$, test de igualdad $p = 0.005$) y en salario por hora ($\hat\tau_F = -0.128$ vs.\ $\hat\tau_I = +0.265$, test $p < 0.001$). Los nulos agregados emergen como resultado mec\'{a}nico de cancelaci\'{o}n entre canales con signos opuestos, no como evidencia de ausencia de efecto. La preregistraci\'{o}n de la formalidad como dimensi\'{o}n central de heterogeneidad queda emp\'{i}ricamente justificada.

Tercero, la lectura cualitativa peruana es consistente con una versi\'{o}n refinada de la \emph{Narrativa D} preregistrada. Las firmas formales que adoptan herramientas con LLM expanden el empleo en ocupaciones expuestas (margen extensivo positivo) y reducen las horas semanales por trabajador (ajuste intensivo). El sector informal experimenta contracci\'{o}n d\'{e}bil del empleo combinada con selecci\'{o}n positiva en salarios entre los trabajadores que permanecen en ocupaciones expuestas. Este cuadro contradice la lectura de complementariedad uniforme del sector formal de la OCDE y entrega un mecanismo distinto: complementariedad sesgada por habilidad, modulada por la segmentaci\'{o}n estructural del mercado laboral.

\paragraph{Limitaciones y trabajo en curso.} Cinco limitaciones del an\'{a}lisis actual orientan la versi\'{o}n siguiente del art\'{i}culo. Primero, el alcance geogr\'{a}fico de la EPEN restringe la representatividad a Lima Metropolitana, no al Per\'{u}. Segundo, la definici\'{o}n de formalidad basada en \texttt{seguro1} es una aproximaci\'{o}n por afiliaci\'{o}n a ESSALUD; la ENAHO anual permite una validaci\'{o}n con la definici\'{o}n est\'{a}ndar OIT basada en pensiones. Tercero, la armonizaci\'{o}n del clasificador ocupacional descansa en pesos posteriores emp\'{i}ricos cuya sensibilidad a esquemas alternativos ser\'{a} contrastada en la escalera de robustez de la \cref{tab:peru_classifier_ladder}. Cuarto, las medidas alternativas de exposici\'{o}n ($\alpha$, $\zeta$, Webb, Felten) y el donut DiD que excluye los trimestres adyacentes al cambio de clasificador permanecen como verificaciones pendientes en la secci\'{o}n de robustez. Quinto, y de orden mayor, los pa\'{i}ses restantes del art\'{i}culo (Costa Rica, Colombia, Ecuador, M\'{e}xico y Uruguay) requieren la replicaci\'{o}n del pipeline para entregar las seis estimaciones $\{\hat\tau^c\}_{c=1}^6$ que constituyen el resultado principal preregistrado.

La versi\'{o}n final del art\'{i}culo contrastar\'{a} los seis pa\'{i}ses lado a lado e identificar\'{a} cu\'{a}l de las cuatro narrativas preregistradas se sostiene en cada pa\'{i}s. Los resultados peruanos sugieren que el agregado regional puede esconder heterogeneidades de signo entre canales formal e informal que merecen estimaci\'{o}n separada en cada pa\'{i}s, y que la evidencia causal sobre el efecto laboral de los LLMs en Am\'{e}rica Latina probablemente difiera de las primeras estimaciones causales de econom\'{i}as de la OCDE no por replicaci\'{o}n del nulo, sino por la mediaci\'{o}n de instituciones de mercado laboral espec\'{i}ficas de la regi\'{o}n.
